% (C) 2026 Ralf Stephan, in collaboration with Claude Code.
% Released under CC0 1.0 Universal (public-domain dedication).
% See https://creativecommons.org/publicdomain/zero/1.0/

\documentclass[11pt]{article}
\usepackage[T1]{fontenc}
\usepackage[utf8]{inputenc}
\usepackage{lmodern}
\usepackage[margin=2.5cm]{geometry}
\usepackage{amsmath,amssymb}
\usepackage{braket}
\usepackage[hidelinks]{hyperref}

% Lean names and paths: verbatim, breakable at . _ /
\DeclareUrlCommand\lean{\urlstyle{tt}}
\def\UrlBreaks{\do\.\do\_\do\/}
% Marks an open problem that blocks checking a main claim of the brief
\newcommand{\blocker}[1]{\textsc{[blocks verification: #1]}~}
% Extra stretch for paragraphs with long, unbreakable Lean names
\emergencystretch=6em

\title{Blocker:\\ Quantum pure states are not enough}
\author{}
\date{}

\begin{document}
\maketitle

\noindent
This file collects everything about the brief issue of \lean{open-problems.tex} that blocks
verification of Theorem~1, together with a proposed resolution and the changes it needs in the
brief. Sections~\ref{sec:blocker} and~\ref{sec:related} are copied from
\lean{open-problems.tex} and \lean{formalization-log.tex}; the rest is new.

\section{The blocker}\label{sec:blocker}
\begin{enumerate}
\item \blocker{brief issue}\textbf{The range of Theorem~1.} The theorem claims consistency for
  $1 \le \dim V < \infty$, while its proof uses $2 \le \dim X$. Making it precise needs four
  choices.
  \begin{itemize}
  \item \emph{Dimension $0$.} The range starts at $1$ if a $0$-dimensional symplectic manifold
    ($n = 0$ degrees of freedom) is allowed, and at $2$ if not. For $\dim V = 1$ the state space
    is a single point, and $P(V)$ has two points exactly when $2 \le \dim V$
    (\lean{subsingleton_projectivization}, \lean{nontrivial_projectivization_iff}).
  \item \emph{Theorem~3 depends on it.} If a point is allowed, Theorem~3 is false for
    $\dim\mathcal H = 1$: the CST ensembles of a point and the density operators of
    $\mathbb C$ are both a single point, and the constant map identifies them
    (\lean{cst_qeDens_consistent_of_finrank_one}). So the brief cannot keep both
    $1 \le \dim V$ in Theorem~1 and Theorem~3 for every $\mathcal H$.
  \item \emph{The empty space.} For $\dim V = 0$, $P(V)$ is empty
    (\lean{isEmpty_projectivization}), which is a symplectic manifold of every dimension. Either
    lower bound, $1$ or $2$, needs CST to exclude it, for instance by requiring an ensemble.
  \item \emph{What ``represented by'' keeps.} If $X$ and $P(V)$ need only be in bijection, the
    ``only if'' is false: $P(\ell^2)$ and the sphere $P(\mathbb C^2)$ both have $\mathfrak c$
    points (\lean{nonempty_projectivization_equiv}), and the sphere's symplectic structure can be
    carried over. The theorem needs $V$ Hausdorff and $X$ to carry the quotient topology of
    $P(V)$; finiteness then comes from CST's Liouville measure $\omega^n$ with finite $n$, not
    from $X$ being symplectic (Log, ``Problems found'').
  \end{itemize}
  Proposed: exclude dimension $0$ and the empty space, so that Theorem~1 reads
  $2 \le \dim V < \infty$ for a Hausdorff $V$ with $X = P(V)$ as topological spaces. Theorem~3
  and its Lean proof then stand as they are. The alternative keeps $1 \le \dim V$, also
  excludes the empty space and restricts Theorem~3 to $\dim\mathcal H \ge 2$. (Log, ``Problems
  found'', first three items.)
\end{enumerate}

\section{Related items}\label{sec:related}
From \lean{open-problems.tex}, ``From the formalization'' (the other blocker, a formalization
gap, which the choices above affect):
\begin{enumerate}
\item \blocker{formalization gap}\textbf{Theorem~1.} Mathlib has no symplectic manifolds and no
  manifold structure on projective space, so neither direction of Theorem~1 is formalized: that
  $P(V)$ with the Fubini--Study form is symplectic for $\dim V < \infty$, and that a symplectic
  $P(V)$ forces $\dim V < \infty$. Theorem~3 for $\dim\mathcal H = \infty$ rests on the second
  direction, so it is formalized only for $\dim\mathcal H < \infty$. The second direction
  needs the topology of $P(V)$; as sets it is false (\lean{nonempty_projectivization_equiv}).
  (Log, ``Results''.)
\end{enumerate}
From \lean{open-problems.tex}, ``From the brief'':
\begin{enumerate}
\setcounter{enumi}{1}
\item \textbf{Ensembles supported on a subspace.} The proof of Theorem~3 compares the ensembles
  with support in a two-dimensional subspace $A$. For $\dim\mathcal H \ge 3$ no CST ensemble has
  support in $P(A)$, which has Liouville measure zero. Whether the brief means the classical
  ensembles of the subsystem $P(A)$, with its own Liouville measure, is open. The formalization
  compares the whole ensemble spaces instead. (Log, ``Problems found'', second item.)
\end{enumerate}
From \lean{formalization-log.tex}, ``Problems found in the brief'':
\begin{itemize}
\item Theorem~1 states $1 \le \dim V$, while its proof uses $2 \le \dim X$. For $\dim V = 1$ the
  state space $P(V)$ is a single point (\lean{subsingleton_projectivization}), a manifold of
  dimension $0$. So either $0$-dimensional symplectic manifolds count, and the proof's bound
  should be $0 \le \dim X$, or they do not, and the theorem should say $2 \le \dim V < \infty$.
  The choice also decides Theorem~3. If a point counts, then for $\dim\mathcal H = 1$, which
  QE-DENS allows, the CST ensembles of the point are $\{\delta\}$ and the density operators are
  $\{1\}$, and the constant map identifies the two (\lean{cst_qeDens_consistent_of_finrank_one}):
  Theorem~3 fails. The Lean proof of Theorem~3 avoids this through \lean{HasManyDisjointSets},
  which fails for a point (\lean{not_hasManyDisjointSets_dirac}). So either Theorem~1 says
  $2 \le \dim V$, or Theorem~3 needs $\dim\mathcal H \ge 2$.
\item For $\dim V = 0$, $P(V)$ is empty (\lean{isEmpty_projectivization}), and the empty set is a
  symplectic manifold of every dimension. Either lower bound, $1$ or $2$, needs CST to exclude it,
  for instance by requiring an ensemble: there is no probability measure on the empty set.
\item The ``only if'' of Theorem~1 needs the identification of $X$ with $P(V)$ to respect the
  topology. As sets, $P(\ell^2)$ and $P(\mathbb C^2)$, the sphere, are in bijection: both have
  $\mathfrak c$ points (\lean{nonempty_projectivization_equiv}). Carrying the symplectic structure
  of the sphere over along such a bijection would make $P(\ell^2)$ a symplectic manifold. The
  theorem holds when $V$ is Hausdorff and $X$ has the quotient topology of $P(V)$; for a
  non-Hausdorff $V$, $P(V)$ need not be Hausdorff even in finite dimension, and the ``if'' fails.
  ``Since $X$ is a symplectic manifold, $\dim X < \infty$'' holds only for finite-dimensional
  symplectic manifolds, and the projective space of an infinite-dimensional Hilbert space is an
  infinite-dimensional one. The finiteness comes from CST itself: the Liouville measure $\omega^n$
  needs a finite $n$.
\item CST says $p(\xi_1\mid\xi_2) = 0$ ``for all $\xi_1, \xi_2 \in X$'', which includes
  $p(\xi\mid\xi) = 0$. It should say $\xi_1 \ne \xi_2$, as the Remark after Example~2
  ($p(\psi\mid\phi) = \delta_{\psi\phi}$) has it.
\end{itemize}

\section{The author's proposal (2026-11-02)}
Paraphrased, from a reply on another brief that the author means for all briefs: the briefs get
a section of \emph{mathematical prerequisites}, which states what Lean should prove and is fully
formalized. The relations between physical conditions, such as Theorems~1 and~3, are argued
informally from these prerequisites and are checked by reading, not by Lean. The author will
rewrite the briefs accordingly.

Here the four choices are not about how to read a physical condition: each one decides which
mathematical statement is true. So they belong in the prerequisites, with their hypotheses
explicit, and in the conditions, which must exclude the degenerate cases.

\section{Proposed resolution}
\paragraph{Guess.} The proposal of Section~\ref{sec:blocker}, with one refinement:
\begin{itemize}
\item CST requires at least one degree of freedom ($n \ge 1$) and a non-empty $X$;
\item QP-RAY takes $V$ to be a complex normed space, for instance a Hilbert space, and
  ``represented by'' to mean that $X$ is homeomorphic to $P(V)$ with the quotient topology;
\item Theorem~1 reads: QP-RAY and CST are consistent if and only if $2 \le \dim V < \infty$;
\item Theorem~3 stands for every $\mathcal H$, including $\dim\mathcal H = 1$, where it holds
  because a point is not a CST state space.
\end{itemize}
The refinement is ``normed'' in place of ``Hausdorff''. For a normed $V$ the infinite-dimensional
case has the short proof of prerequisite P3 below, which uses continuous linear functionals. A
Hausdorff topological vector space may have none (for instance $L^p$ with $0 < p < 1$), and then
that proof does not apply; whether the claim of the log still holds there is not checked.

\paragraph{Why.}
\begin{itemize}
\item \emph{A system with no degrees of freedom is not a physical system.} A phase space that
  is a single point has nothing to prepare, measure or evolve, and neither has a quantum system
  with one state. Physicists do not count either. Excluding them costs nothing.
\item \emph{Theorem~3 is the point of the brief.} With $n \ge 1$, Theorem~3 holds for every
  QE-DENS system without exception, and its Lean proof stands as it is. Keeping
  $1 \le \dim V$ would only buy a trivial exception to the brief's main result.
\item \emph{Every system has a state.} The empty space has no states and no ensembles, so it
  is not a physical system either. Requiring $X \ne \emptyset$, or an ensemble, is a
  condition no one would dispute.
\item \emph{Topology is physical.} In this programme the topology of a state space comes from
  what can be verified experimentally: states that are close cannot be told apart by finite
  measurements. A bare bijection can scramble arbitrarily close states, so it is not a
  physical notion of ``represented by''. The natural topology on $P(V)$ is the quotient
  topology, which for a Hilbert space agrees with the Fubini--Study metric.
\item \emph{Normed spaces suffice.} QE-DENS uses Hilbert spaces, and every space used in
  quantum mechanics is at least normed. Restricting QP-RAY to normed spaces loses no physics.
\end{itemize}

\paragraph{Confidence.}
\begin{itemize}
\item High: $2 \le \dim V < \infty$, excluding the empty space, and the topological reading of
  ``represented by''. Each excludes a degenerate case or follows a standard operational
  principle, and together they keep Theorem~3 without exception.
\item Medium-high: ``normed'' in place of ``Hausdorff''. It is the safe choice; the author may
  prefer the general topological vector space of QP-RAY, which then needs its own proof.
\end{itemize}

\section{Mathematical prerequisites to add to the brief}
Each item is a purely mathematical statement for Lean.
\begin{enumerate}
\item[P1.] \emph{Size of $P(V)$.} For a complex vector space $V$: $P(V)$ is empty if and only if
  $\dim V = 0$, is a single point if $\dim V = 1$, and has at least two points if and only if
  $\dim V \ge 2$.\par{\raggedright
  Lean: \lean{isEmpty_projectivization}, \lean{nonempty_projectivization_iff},
  \lean{subsingleton_projectivization}, \lean{nontrivial_projectivization_iff}.\par}
\item[P2.] \emph{Finite dimension.} If $\dim V = k + 1$ with $1 \le k < \infty$, then $P(V)$
  with the Fubini--Study form is a compact symplectic manifold of real dimension $2k$.\\
  Lean: not formalized; Mathlib has no symplectic manifolds and no manifold structure on
  $P(V)$.
\item[P3.] \emph{Infinite dimension.} If $V$ is a complex normed space of infinite dimension,
  then $P(V)$ with the quotient topology is not locally compact, hence not homeomorphic to a
  finite-dimensional manifold. Proof: for $[v] \in P(V)$ choose a continuous linear functional
  $f$ with $f(v) = 1$. The set $\{[w] : f(w) \ne 0\}$ is open in $P(V)$ and homeomorphic to the
  affine hyperplane $f^{-1}(1)$, an infinite-dimensional normed affine space, which is not
  locally compact by Riesz's lemma.\\
  Lean: not formalized. It needs no symplectic geometry; Mathlib has the Hahn--Banach and Riesz
  results.
\item[P4.] \emph{Theorem~3, the mathematical core.} Let $\mu$ be a measure with, for every $N$,
  $N$ disjoint sets of positive finite measure (for instance the Lebesgue measure on a non-empty
  open set of $\mathbb R^{2n}$, $n \ge 1$, the Liouville measure of a Darboux chart), and let
  $\mathcal H$ be a finite-dimensional Hilbert space. No injective map from the probability
  measures absolutely continuous with respect to $\mu$ to the density operators of
  $\mathcal H$ preserves mixtures.\\
  Lean: \lean{cst_qeDens_inconsistent}, \lean{hasManyDisjointSets_restrict}. Not formalized:
  the passage from a Darboux chart to the whole manifold.
\end{enumerate}
Theorem~3 for $\dim\mathcal H = \infty$ then follows informally: by P3, $P(\mathcal H)$ is not a
finite-dimensional manifold, so it is not a CST state space.

\section{Other changes to the brief}
\begin{enumerate}
\item \emph{CST.} Require $n \ge 1$ and $X \ne \emptyset$. Replace ``for all
  $\xi_1, \xi_2 \in X$'' by ``for all $\xi_1 \ne \xi_2$ in $X$''.
\item \emph{QP-RAY.} Take $V$ to be a complex normed space, or say which topological vector
  spaces are meant, and say that $X$ is homeomorphic to $P(V)$ with the quotient topology.
\item \emph{Theorem~1.} Change ``$1 \le \dim(V) < \infty$'' to ``$2 \le \dim(V) < \infty$''.
\item \emph{Proof of Theorem~1.} Replace ``Since $X$ is a symplectic manifold, then
  $2 \le \dim(X) < \infty$'' by: CST's Liouville measure $\omega^n$ needs a finite $n \ge 1$,
  so $X$ is a manifold of finite dimension $2n \ge 2$; then $\dim V < \infty$ by P3 and
  $\dim V \ge 2$ by P1, and conversely P2 gives the ``if''.
\item \emph{Proof of Theorem~3.} Compare the whole ensemble spaces, as in P4, instead of the
  ensembles with support in a two-dimensional subspace $A$, which do not exist for
  $\dim\mathcal H \ge 3$ (item~2 of \lean{open-problems.tex}). Cite P3 for
  $\dim\mathcal H = \infty$.
\end{enumerate}

\paragraph{When the blocker closes.} The brief issue closes once the brief makes these choices.
The formalization gap of Section~\ref{sec:related} remains until P2 is formalized; P3 can be
formalized now and would close the infinite-dimensional case of Theorem~3.

\end{document}
